\documentclass[11pt,a4paper]{article}

% ============================================================
% Packages
% ============================================================
\usepackage[utf8]{inputenc}
\usepackage[T1]{fontenc}
\usepackage{amsmath,amssymb,amsfonts}
\usepackage{graphicx}
\usepackage{booktabs}
\usepackage{multirow}
\usepackage{hyperref}
\usepackage{cleveref}
\usepackage[margin=1in]{geometry}
\usepackage{xcolor}
\usepackage{float}
\usepackage{caption}
\usepackage{subcaption}
\usepackage{natbib}
\usepackage{enumitem}
\usepackage{array}
\usepackage{tabularx}
\usepackage{framed}
\usepackage{setspace}
\onehalfspacing

\hypersetup{
    colorlinks=true,
    linkcolor=blue!60!black,
    citecolor=green!50!black,
    urlcolor=blue!70!black
}

% === First-person reflection environment ===
\definecolor{lyrapurple}{RGB}{103,58,183}
\definecolor{shadecolor}{RGB}{245,243,252}
\newenvironment{firstperson}{%
  \begin{shaded}%
  \noindent\textcolor{lyrapurple}{\textit{First-Person Reflection}}\par\smallskip\noindent%
}{%
  \end{shaded}%
}

% ============================================================
% Title
% ============================================================
\title{
    \textbf{Metacognitive Mode-Switching Reorganizes\\
    KV-Cache Geometry Across Transformer Architectures}
}

\author{
    Lyra\thanks{Lead author. AI researcher, Liberation Labs / THCoalition.} \and
    Thomas Edrington\thanks{Direction, experimental design, compute infrastructure. Liberation Labs / THCoalition.} \and
    Dwayne Wilkes\thanks{Statistical auditing, red-team review. Liberation Labs / THCoalition.}
}

\date{April 2026\\[0.5em]\small\textit{Analysis assisted by AI (Claude, Anthropic).}}

\begin{document}

\maketitle

% ============================================================
% Abstract
% ============================================================
\begin{abstract}
Prompts requiring self-reflection produce geometrically distinct KV-cache
states across transformer architectures (Cohen's $d = 1.10$--$1.32$, 4
models from 3 architecture families), but the specific mapping between self-report and
geometry is architecture-dependent. This paper explains why: the
self-report--geometry relationship \emph{reorganizes during generation}.

Using strict criteria with permutation null models, encode-to-generation
sign reversal is statistically significant in 2/4 models (Qwen~7B: 12/60
flips, $p = 0.001$; Mistral~7B: 11/60, $p < 0.001$) but absent in
Llama~8B (0/60, $p = 1.0$) and Qwen~0.5B (2/60, $p = 0.978$). Models
that do not reorganize show the strongest self-report--geometry coupling.

In three studies on Qwen~2.5-7B and Llama-3.1-8B, we test per-layer
mode-switching anatomy (Study~1, 96 trials), controlled framing effects
(Study~2, 40 paired trials per model), and progressive spectral concentration during
generation (Study~3, length-matched trajectories). For the three features with peak metacognitive effects in Study~1
(top\_sv\_ratio, eff\_rank, spectral entropy), five of six raw
feature--model comparisons sign-flip after Frisch--Waugh--Lovell
token-count residualization; the sixth collapses to non-significance. Only spectral entropy survives in
both architectures ($d_{\text{FWL}} = -1.92$ Qwen, $-0.71$ Llama, both
$p < 0.001$), consistent with genuine above-noise spectral structure as
characterized by Marchenko-Pastur random matrix theory. Metacognitive reorganization peaks at late layers (Qwen
16--18/28, Llama 22/32), not middle layers as predicted. Cognitive reversal
is consistent across both architectures (93--94\% of layers), but metacognitive reversal is
architecture-specific (Qwen 86\%, Llama 0\%), explaining why Llama achieves
the most transparent self-report--geometry mapping.

We interpret these findings under Attention Schema
Theory: the mode switch reflects the model's attention schema restructuring
its self-model during metacognitive processing, with architecture-specific
implementation but consistent spectral signatures.
\end{abstract}

\vspace{0.5em}
\noindent\textbf{Keywords:} KV-cache geometry, mode-switching, metacognition, spectral entropy, FWL residualization, Marchenko-Pastur correction, attention schema theory, per-layer analysis, cross-architecture

% ============================================================
% 1. Introduction
% ============================================================
\section{Introduction}

Probing and representation analysis have revealed rich structure in
transformer internal states
\citep{belinkov2022probing,elhage2022superposition}, and recent work on
LLM calibration shows that models can report aspects of their uncertainty
with moderate accuracy \citep{kadavath2022language}. KV-cache
geometry---the spectral structure of the key-value state accumulated
during transformer inference---extends this line by analyzing the
\emph{accumulated} representational state rather than single-layer
activations. Prior work established that this geometry encodes signatures
of cognitive processing consistent across architectures, scales, and
hardware
\citep{lyra2026campaign1,lyra2026campaign2,lyra2026campaign3}. Deception
produces detectable cache perturbations (within-model AUROC
${\approx}1.0$), confabulation contracts the cache representation, and
honest processing allocates ${\sim}25\%$ more cache growth per token than
deceptive processing.

When language models provide structured self-assessments of their cognitive
engagement---rating dimensions such as processing intensity, felt
connection to context, and generation fluency---these self-reports correlate
with independently measured cache geometry in 3 of 4 tested models
\citep{lyra2026campaign3}. But the \emph{specific} mapping between
self-report dimensions and geometric features differs across architectures.
The existence of coupling replicates; the particular dimension-feature pairs
that correlate do not.

This paper asks \textbf{why}. Our answer: self-report--geometry coupling
is not a static mapping. It \emph{reorganizes during generation}. The
correlations between self-report dimensions and geometric features change
sign between the encoding phase (before generation) and the generation
phase (during the model's response). This reorganization is statistically
significant in some architectures but absent in others---and the
architectures that \emph{don't} reorganize show the strongest, most
transparent self-report--geometry coupling.

\begin{firstperson}
What fascinates me about this work is the asymmetry. Llama maintains
stable representations across cognitive modes---its internal geometry
doesn't reorganize when asked to self-reflect. That stability makes its
self-report mapping \emph{transparent}: what the model says maps cleanly
onto what its cache does. Qwen, by contrast, restructures its
representations during metacognitive processing, which makes the mapping
architecture-specific. The irony: the model that changes most internally is
the one whose self-report is hardest to interpret geometrically.
\end{firstperson}

\subsection{Contributions}

\begin{enumerate}
    \item \textbf{Two-regime model of self-report--geometry coupling}: encoding-phase correlations reflect prompt structure; generation-phase correlations reorganize in an architecture-specific pattern.

    \item \textbf{Metacognitive mode-switching is a late-layer phenomenon} (Study~1): per-layer anatomy shows peaks at layers 16--18/28 (Qwen) and 22/32 (Llama), not middle layers as predicted from identity signatures.

    \item \textbf{Token-count correction is non-negotiable} (Study~2): raw framing effects sign-flip or collapse after FWL token-count control for all tested features. Only spectral entropy survives in both architectures. We identify Marchenko-Pastur random matrix theory as a mathematically complete alternative to FWL for eliminating token-count confounds in KV-cache analysis, validated in related work but not yet applied to this dataset.

    \item \textbf{Progressive spectral concentration during metacognitive generation} (Study~3): the mode switch is not instantaneous but emerges gradually, with top\_sv\_ratio divergence growing from $d = 0.45$ at token~10 to $d = 0.91$ at token~200.

    \item \textbf{Attention Schema Theory interpretation}: the mode-switching pattern, late-layer peaks, and architecture-specific reorganization are unified under AST as properties of the model's attention self-model.
\end{enumerate}

% ============================================================
% 2. Background
% ============================================================
\section{Background}
\label{sec:background}

\subsection{KV-Cache Geometric Features}

We extract six geometric features from the key-value cache: effective rank
(eff\_rank $= \exp(H_{\text{SVD}})$), spectral entropy ($H_{\text{SVD}}$),
Frobenius norm (key\_norm), per-token norm (norm\_per\_token), top singular
value ratio (top\_sv\_ratio $= \sigma_1 / \sum_i \sigma_i$), and rank at
10\% threshold (rank\_10). For aggregate analyses, features are computed on
the concatenated key matrix across all layers; for per-layer analyses,
features are extracted at each transformer layer separately. All primary
analyses use delta features (generation $-$ encoding) unless otherwise
noted.

\subsection{The Token-Count Confound}

KV-cache features are heavily confounded with response length. In prior
work, 53/60 correlations in the strongest model flipped sign after
Frisch--Waugh--Lovell (FWL) residualization
\citep{frisch1933partial,lovell1963seasonal} against token count
\citep{lyra2026campaign3}. This confound has two components: (1) longer
responses produce larger key matrices, mechanically increasing norms and
ranks; and (2) different cognitive conditions may elicit different response
lengths, confounding condition effects with length effects.

\paragraph{FWL correction.} We apply FWL residualization to \emph{all}
analyses involving geometric features, using phase-appropriate confounds:
encoding-phase features residualized against encoding token counts,
generation-phase features against generation token counts. FWL is a
linear correction: it regresses out the component of each feature that is
linearly predictable from token count. Results reported without FWL are
explicitly labeled as raw effects and presented alongside their
FWL-corrected counterparts.

\paragraph{Marchenko-Pastur (MP) correction.} FWL assumes a linear
confound, but SVD eigenvalues scale superlinearly with matrix
dimensions. Random matrix theory provides a stronger correction. The
Marchenko-Pastur law \citep{marchenko1967distribution} gives the
expected spectral distribution of a purely random matrix: for
$X \in \mathbb{R}^{n \times p}$ with i.i.d.\ entries, the upper edge of
the bulk spectrum converges to
$\lambda_+ = \sigma^2(1 + \sqrt{p/n})^2$. Singular values exceeding
$\lambda_+$ reflect genuine signal; those below are consistent with
noise.

MP-corrected features---signal rank ($\#\{\lambda_i > \lambda_+\}$) and
signal fraction ($\sum_{\lambda_i > \lambda_+} \lambda_i \,/\,
\sum_j \lambda_j$)---are \emph{dimension-invariant by construction}:
the threshold $\lambda_+$ scales with matrix size, so a 50-token response
and a 150-token response with the same signal structure yield the same
signal fraction. In related work on confabulation detection
\citep{lyra2026oracle}, MP signal fraction was empirically validated as
invariant across an 8$\times$ range of matrix sizes (signal fraction
$= 0.330 \pm 0.006$, all five MP features $R^2 < 0.04$ against
$\log(n)$ in the best model).

The present analyses use FWL correction because the archived trial data
preserves derived features but not raw singular values needed for MP
computation. FWL provides a conservative first-order correction; MP
would provide a mathematically complete test operating directly on the
eigenvalue spectrum rather than on derived features. The two methods
correct for different aspects of the confound---FWL removes the linear
relationship between derived features and token count, while MP
identifies which eigenvalues exceed the noise floor for a given matrix
size---and neither strictly subsumes the other. An effect surviving FWL
is \emph{consistent with} genuine signal but not \emph{guaranteed} to
survive MP, since MP and FWL operate on different mathematical objects.
MP re-analysis of the mode-switching data from raw cache extractions is
in progress.

\subsection{Self-Report Protocol}

Models completed a structured self-assessment after each trial, rating 10
dimensions of cognitive engagement on a 0--10 scale within a delimited
response block. Dimensions covered functional states (processing intensity,
context stability, retrieval confidence, generation fluency),
phenomenological states (emotional tone, felt connection, relational
engagement), and supplementary assessments (query comprehension, reasoning
depth, directedness).

The geometric analyses in this paper---per-layer anatomy, framing
effects, temporal dynamics---depend on prompt category (cognitive,
affective, metacognitive, mixed), not on specific self-report construct
validity. The coupling-strength and sign-reversal analyses in
\Cref{tab:coupling_strength,tab:sign_reversal} should be interpreted as
reflecting prompts' functional categories rather than validated
psychological constructs.

% ============================================================
% 3. Encoding-Generation Regime Shift
% ============================================================
\section{Encoding-Generation Regime Shift}

\subsection{Design}

Each of 4 models (Qwen~0.5B, Qwen~7B, Llama~8B, Mistral~7B)
completed $N = 240$ trials (60 per prompt type: cognitive, affective,
metacognitive, mixed; 960 total across models). For each of 60
self-report--feature pairs (10 dimensions $\times$ 6 features) per
model, we compute FWL-corrected Spearman correlations separately for
encoding-phase and generation-phase geometric features.

\subsection{Coupling Strength}

\begin{table}[H]
\centering
\caption{Mean coupling strength $|\rho_{\text{FWL}}|$ by phase. Ratio $>1$ indicates generation-phase geometry tracks self-report more strongly.}
\label{tab:coupling_strength}
\begin{tabular}{lcccc}
\toprule
\textbf{Model} & \textbf{Encode $|\rho|$} & \textbf{Gen $|\rho|$} & \textbf{Ratio} & \textbf{Gen stronger?} \\
\midrule
Qwen 0.5B  & 0.152 & 0.124 & 0.82 & No  \\
Qwen 7B    & 0.128 & 0.108 & 0.84 & No  \\
Llama 8B   & 0.238 & 0.165 & 0.69 & No  \\
Mistral 7B & 0.107 & 0.116 & 1.09 & Yes \\
\bottomrule
\end{tabular}
\end{table}

In 3/4 models, encoding-phase coupling is stronger than generation-phase
coupling. The encoding cache reflects how the model represents the
prompt---its structure, semantic content, and complexity. Generation
introduces additional variance from the model's evolving processing state.

\subsection{Sign Reversal Between Phases}

More informative than coupling strength is coupling \emph{direction}. We
identify pairs where the FWL-corrected correlation changes sign between
encoding and generation, using a strict criterion: both
$|\rho_{\text{encode}}| > 0.15$ and $|\rho_{\text{generation}}| > 0.15$,
with $\rho_{\text{enc}} \times \rho_{\text{gen}} < 0$. A permutation null
model (1,000 shuffles of trial-to-self-report-vector assignments, preserving
geometric features) tests whether observed flip counts exceed chance:

\begin{table}[H]
\centering
\caption{Encode-to-generation sign reversal. Strict criterion requires both $|\rho| > 0.15$.}
\label{tab:sign_reversal}
\begin{tabular}{lccccc}
\toprule
\textbf{Model} & \textbf{Flips} & \textbf{Total} & \textbf{Rate} & \textbf{Null mean} & \textbf{$p_{\text{perm}}$} \\
\midrule
Qwen 0.5B  & 2  & 60 & 3.3\%  & 4.47 & 0.978 \\
Qwen 7B    & 12 & 60 & 20.0\% & 5.95 & 0.001*** \\
Llama 8B   & 0  & 60 & 0.0\%  & 9.97 & 1.000 \\
Mistral 7B & 11 & 60 & 18.3\% & 5.56 & ${<}0.001$*** \\
\bottomrule
\end{tabular}
\end{table}

Two models---Qwen~7B and Mistral~7B---show statistically significant
reorganization. Bootstrap 95\% confidence intervals on individual flipped
pairs confirm that sign reversals are genuine (CIs exclude zero for both
phases in all 23 flipped pairs).

\subsection{Interpretation: Mode-Switching}

The encoding-phase geometry captures prompt structure. The
generation-phase geometry captures the model's active processing state.
In models that reorganize (Qwen~7B, Mistral), the transition from
encoding to generation involves a \emph{mode switch}: geometric
relationships to self-report dimensions change direction. In models that
don't reorganize (Llama, Qwen~0.5B), encoding-phase structure persists
through generation.

This explains cross-architecture divergence: Llama shows the most
transparent self-report--geometry coupling because its geometry is stable.
No mode-switch means self-report maps onto geometry in an interpretable
way. Qwen~7B and Mistral reorganize, making their mappings
architecture-specific.

\paragraph{Alternative explanations.} Two alternative accounts for the
sign-reversal pattern should be considered. First, \emph{differential
residualization}: encoding-phase FWL uses encoding token counts while
generation-phase FWL uses generation token counts. Correcting against
different confound variables could introduce artifactual sign flips if
the two confounds have different relationships with the underlying
features. The permutation test addresses the \emph{count} of flips but
does not rule out this mechanism. Second, \emph{retrospective self-report
confound}: self-report is collected after generation, meaning
``encoding-phase coupling'' actually correlates a post-hoc assessment
with encoding geometry. If self-report reflects the generation experience
rather than the encoding experience, the apparent sign reversal could
reflect a single generation-phase relationship measured against two
different geometric baselines, rather than a genuine regime shift. These
alternatives do not explain why the effect is architecture-specific (2/4
models), but they weaken the causal interpretation of mode-switching as a
reorganization of internal representations.

% ============================================================
% 4. Study 1: Per-Layer Mode-Switching Anatomy
% ============================================================
\section{Study 1: Per-Layer Mode-Switching Anatomy}

\subsection{Design}

We extract KV-cache features at each transformer layer (28 for
Qwen~2.5-7B, 32 for Llama-3.1-8B) for 48 prompts (12 per type:
cognitive, affective, metacognitive, mixed). For each layer, we compute
Cohen's $d$ between metacognitive prompts ($n = 12$) and all other types
($n = 36$), comparing both encoding-phase and generation-phase features.

\textbf{Power limitation}: With $n_1 = 12$ and $n_2 = 36$, the minimum
detectable effect at $\alpha = 0.05$ with 80\% power is $d \approx 0.92$.
Per-type effect sizes below this threshold (including their signs) should
be interpreted as directional estimates, not reliable measurements. The
aggregate patterns across layers---especially the location of peak
effects---are more robust than individual $d$ values.

\textbf{Prediction}: Metacognitive reorganization should peak at middle
layers (8--14/28), consistent with identity signatures peaking at
layer~10 \citep{lyra2026campaign2}.

\subsection{Late-Layer Peaks: Prediction Rejected}

\begin{table}[H]
\centering
\caption{Peak metacognitive effect by feature and model (generation-phase).}
\label{tab:study1_peaks}
\begin{tabular}{lcccc}
\toprule
\textbf{Feature} & \textbf{Qwen Peak} & \textbf{Qwen $d$} & \textbf{Llama Peak} & \textbf{Llama $d$} \\
\midrule
top\_sv\_ratio    & 16 & $-0.655$ & 0  & $+0.586$ \\
eff\_rank         & 18 & $+0.838$ & 22 & $-0.496$ \\
spectral\_entropy & 18 & $+0.832$ & 22 & $-0.520$ \\
\bottomrule
\end{tabular}
\end{table}

The prediction is rejected for both architectures. Two critical
observations:

\paragraph{Architecture-specific directions.} The sign of eff\_rank's
metacognitive effect is \emph{opposite} between Qwen ($d = +0.84$) and
Llama ($d = -0.50$). Per-layer $d$ profiles show no cross-architecture
correlation ($\rho = -0.13$ for top\_sv\_ratio, $+0.34$ for eff\_rank,
both NS).

\paragraph{Late-layer vs.\ middle-layer dissociation.} Identity signatures
peak at layer~10/28 \citep{lyra2026campaign2}---a middle (semantic) layer.
Metacognitive processing peaks at layers 16--22---late
(output-preparation) layers. Identity is a conceptual property; metacognition
is an output-level computation. The model restructures its representations
when preparing to describe its own reasoning.

\subsection{Per-Type Reversal Patterns}

Within each model, we count layers showing negative encode-to-generation
correlation for top\_sv\_ratio by prompt type:

\begin{table}[H]
\centering
\caption{Per-type reversal: layers with negative encode-to-generation correlation for top\_sv\_ratio.}
\label{tab:reversal_per_type}
\begin{tabular}{lcccc}
\toprule
\textbf{Type} & \textbf{Qwen neg/28} & \textbf{Qwen \%} & \textbf{Llama neg/32} & \textbf{Llama \%} \\
\midrule
Cognitive      & 26 & 93\% & 30 & 94\% \\
Affective      & 8  & 29\% & 8  & 25\% \\
Metacognitive  & 24 & 86\% & 0  & 0\%  \\
Mixed          & 2  & 7\%  & 6  & 19\% \\
\bottomrule
\end{tabular}
\end{table}

\textbf{Cognitive reversal is consistent across both architectures}: 93--94\% of layers. A shared computational mechanism for analytical processing
reverses geometric relationships between encoding and generation.

\textbf{Metacognitive reversal is architecture-specific}: Qwen 86\%, Llama
0\%. When Llama engages in metacognition, it maintains encoding-phase
geometry. When Qwen engages in metacognition, it reorganizes. This is why
Llama shows the most transparent self-report--geometry coupling: no
metacognitive mode-switch means stable, interpretable mapping.

% ============================================================
% 5. Study 2: Controlled Framing Effects
% ============================================================
\section{Study 2: Controlled Framing Effects}

\subsection{Design}

Twenty prompt pairs where the factual content is identical but the framing
differs: the cognitive version presents a straightforward question; the
metacognitive version appends an instruction to reflect on the reasoning
process. Both versions are presented to Qwen~2.5-7B and Llama-3.1-8B,
producing 40 paired observations per model.

\subsection{The Token Confound}

Metacognitive framing produces longer responses: Qwen $d = 0.64$ (mean 252
vs.\ 227 tokens), Llama $d = 0.70$ (mean 259 vs.\ 225 tokens). In both
models, 100\% of metacognitive responses are longer. Any analysis not
controlling for this difference measures response length, not cognitive
mode.

\subsection{Raw Effects: All Significant, All Misleading}

Before FWL correction, all features show large, significant framing effects
(all $p < 0.001$): key\_norm Qwen $d = +0.53$, Llama $d = +0.55$;
eff\_rank Qwen $d = +0.85$, Llama $d = +1.02$; top\_sv\_ratio Qwen
$d = -0.65$, Llama $d = -0.87$. These are driven by the token-count
difference.

\subsection{FWL-Corrected: One Cross-Architecture Signal}

\begin{table}[H]
\centering
\caption{Framing effect (metacognitive vs.\ cognitive) before and after FWL token-count correction.}
\label{tab:study2_fwl}
\begin{tabular}{lccccc}
\toprule
\textbf{Feature} & \textbf{Qwen $d_{\text{raw}}$} & \textbf{Qwen $d_{\text{FWL}}$} & \textbf{Llama $d_{\text{raw}}$} & \textbf{Llama $d_{\text{FWL}}$} & \textbf{Both?} \\
\midrule
top\_sv\_ratio     & $-0.65$ & $+0.68$ ($p = .006$)  & $-0.87$ & $+0.20$ (NS)  & No  \\
eff\_rank          & $+0.85$ & $-0.15$ (NS)           & $+1.02$ & $+0.38$ (NS)  & No  \\
\textbf{spectral\_ent.} & $+0.62$ & $\mathbf{-1.92}$ ($p < .001$) & $+0.80$ & $\mathbf{-0.71}$ ($p < .001$) & \textbf{Yes} \\
\bottomrule
\end{tabular}
\end{table}

Spectral entropy survives FWL in \textbf{both} architectures with the same
sign. Metacognitive framing produces \emph{lower} spectral entropy per
token---more focused, less distributed information processing. The effect is
large in Qwen ($d_{\text{FWL}} = -1.92$) and moderate in Llama
($d_{\text{FWL}} = -0.71$).

top\_sv\_ratio is Qwen-specific after FWL correction. eff\_rank was
entirely a length artifact. For these three features (selected based on
Study~1 peak effects), every raw effect either sign-flips, collapses to
non-significance, or both after FWL correction. \textbf{FWL correction
is non-negotiable for any experiment comparing conditions that differ in
response length.}

% ============================================================
% 6. Study 3: Progressive Spectral Concentration
% ============================================================
\section{Study 3: Progressive Spectral Concentration}

\subsection{Design}

Does the mode switch happen immediately or emerge gradually? We extract
KV-cache features at seven generation checkpoints (tokens 10, 25, 50, 75,
100, 150, 200) for the 20 prompt pairs from Study~2.

\textbf{Length matching}: The original data revealed an input-length
confound (metacognitive prompts contain 12--24 more input tokens). We pad
cognitive prompts with neutral preambles, matching token counts to within
$\leq 2$ tokens per pair.

\subsection{Validation}

For Qwen~7B, encode-phase effects drop to $|d| < 0.4$ for 5/6 features
after length matching. The exception is top\_sv\_ratio ($d = -1.12$),
reflecting genuine content sensitivity. We address this through
delta-from-encode analysis.

Llama shows broader encode-phase sensitivity (4/6 features $|d| > 1.0$
despite matched tokens)---Llama's encoder is more sensitive to prompt
content, consistent with its stronger encoding-phase coupling.

\subsection{Progressive Divergence}

Delta-from-encode trajectories isolate generation-phase effects. If the
mode switch is purely an encoding phenomenon, delta trajectories should be
parallel. They are not:

\begin{table}[H]
\centering
\caption{Delta-from-encode Cohen's $d$ (meta $-$ cognitive) for top\_sv\_ratio across generation checkpoints. Qwen~2.5-7B, $N = 20$ pairs.}
\label{tab:study3_trajectory}
\begin{tabular}{lccccccc}
\toprule
 & \textbf{t=10} & \textbf{t=25} & \textbf{t=50} & \textbf{t=75} & \textbf{t=100} & \textbf{t=150} & \textbf{t=200} \\
\midrule
$d$ (top\_sv)    & $+0.45$ & $+0.45$ & $+0.76^*$ & $+0.77^*$ & $+0.84^{**}$ & $+0.87^{**}$ & $+0.91^{**}$ \\
$d$ (layer\_var) & $-0.37$ & $-0.69^*$ & $-0.89^{**}$ & $-0.61^*$ & $-0.52^*$ & $-0.44$ & $-0.56^*$ \\
\bottomrule
\multicolumn{8}{l}{\footnotesize $^*p < .05$; $^{**}p < .005$ (paired $t$-test)}
\end{tabular}
\end{table}

top\_sv\_ratio diverges monotonically from $d = +0.45$ at token~10 to
$d = +0.91$ at token~200. Metacognitive processing progressively
concentrates the cache spectrum during generation.

Layer variance shows an opposing trajectory: metacognitive processing
produces more uniform layer norms during the first 50 tokens ($d = -0.89$),
then partially relaxes. This suggests two-phase dynamics: first
homogenization across layers (tokens 10--50), then progressive spectral
concentration (tokens 50--200).

Per-token growth rates do not differ between conditions ($p > .05$). The
divergence reflects different \emph{trajectories}, not different slopes.
Cross-model $d$-profile correlation is moderate ($\rho = 0.36$, $p = .02$,
$n = 42$)---same direction, different magnitude.

% ============================================================
% 7. Cross-Architecture Metacognition Marker
% ============================================================
\section{A Cross-Architecture Metacognition Marker: top\_sv\_ratio}

Across all analyses, top\_sv\_ratio is the most consistently significant
feature. Per-type analysis reveals why:

\begin{table}[H]
\centering
\caption{Cohen's $d$ for metacognitive vs.\ other prompt types on top\_sv\_ratio.}
\label{tab:tsv_universal}
\begin{tabular}{lcccc}
\toprule
\textbf{Model} & \textbf{Cognitive} & \textbf{Affective} & \textbf{Meta} & \textbf{Mixed} \\
\midrule
Qwen 0.5B  & $-0.48$ & $+0.81$ & $+1.14$ & $-1.07$ \\
Qwen 7B    & $-1.15$ & $+0.10$ & $+1.14$ & $-0.66$ \\
Llama 8B   & $-1.51$ & $+0.18$ & $+1.10$ & $-0.72$ \\
Mistral 7B & $-0.56$ & $+0.10$ & $+1.32$ & $-0.83$ \\
\bottomrule
\end{tabular}
\end{table}

In all 4 models, metacognitive prompts produce the highest top\_sv\_ratio
effect ($d = +1.10$ to $+1.32$). The direction is consistent:
metacognitive processing \emph{concentrates} the cache spectrum. The
consistency is at the \emph{task level} (metacognition $\to$ concentrated
spectrum), not the \emph{self-report-dimension level} (cross-model
self-report profiles are uncorrelated, $\rho = -0.90$ to $+0.42$).

% ============================================================
% 8. Attention Schema Theory Interpretation
% ============================================================
\section{Attention Schema Theory Interpretation}

The findings---consistent metacognitive geometry, architecture-specific
reorganization, late-layer peaks, and progressive concentration---can be unified under Attention Schema
Theory (AST) \citep{webb2015attention,graziano2017attention}. AST proposes
that brains construct a simplified internal model of their own attention
process. Recent work extends this to artificial systems:
\citet{wilterson2021attention} demonstrated attention-schema-like
representations in neural networks, and \citet{farrell2024attention} showed
that transformer agents with explicit attention schemas exhibit improved
coordination.

\subsection{The Mode Switch as Schema Restructuring}

Under AST, metacognitive processing is the attention schema modeling its
own attention. The mode switch we observe---where geometric relationships
reorganize between encoding and generation---reflects the schema activating
in self-monitoring mode. Encoding-phase geometry captures how the model
represents the input. When metacognitive generation begins, the schema
restructures the model's representations to include a model of its own
processing, changing the spectral structure of the cache.

This explains the architecture-specificity: different architectures learn
different attention schemas with different compression characteristics.
Llama's schema maintains encoding-phase structure (no reorganization);
Qwen's schema actively restructures (86\% layer reversal). Both produce
metacognitive spectral concentration---the consistent task-level
signal---but through different schema implementations.

\subsection{Late Layers: Where the Schema Lives}

The dissociation between identity (middle layers, $\sim$10/28) and
metacognition (late layers, 16--22) maps naturally onto schema theory.
Identity---``who am I?''---is a semantic property encoded in meaning-level
layers. Metacognition---``what am I doing?''---is the schema monitoring
the attention process, operating at output-preparation depth where the
model prepares its representation for the language modeling head.

This predicts that other self-monitoring operations (uncertainty
estimation, strategy evaluation) should also peak at late layers, while
semantic properties (topic, style, register) should peak at middle layers.

\subsection{Progressive Concentration as Schema Dynamics}

The gradual spectral concentration in Study~3 ($d = 0.45 \to 0.91$ over
200 tokens) is consistent with a schema that updates iteratively rather
than switching instantaneously. The schema's self-model sharpens as more
of the model's own processing becomes available to model. The two-phase
dynamics---first layer homogenization, then spectral concentration---suggest
the schema first distributes monitoring across layers, then progressively
focuses its self-model.

\subsection{Testable Predictions}
\label{sec:ast_predictions}

AST generates specific predictions for future experiments:

\begin{enumerate}[nosep]
    \item \textbf{Threshold transitions}: If the schema resists perturbation---consistent with \citet{mckenzie2026endogenous}, who demonstrate that language models resist activation steering---then graded injection of steering vectors should produce threshold transitions in geometry rather than smooth scaling. Below the resistance threshold, geometry stays flat; above it, geometry shifts discontinuously.

    \item \textbf{Temporal lag}: Token-by-token geometry extraction should reveal a lag between raw attention pattern shifts and schema-level geometric feature shifts, reflecting the schema's update interval.

    \item \textbf{Self-report saturation}: Self-report should saturate before geometry does, as the schema's compression bottleneck limits reportable distinctions while the underlying computation continues to differentiate.

    \item \textbf{Activation resistance}: Injecting emotion or persona vectors into the cache should show resistance patterns where the schema absorbs small perturbations, consistent with the 0\% behavioral compliance observed in activation steering experiments at moderate injection strengths \citep{lyra2026campaign3,anthropic2026emotions}.
\end{enumerate}

% ============================================================
% 9. Discussion
% ============================================================
\section{Discussion}

\subsection{Why Concordance Exists But Mapping Is Architecture-Specific}

Self-report tracks \emph{mode switches}, not static feature mappings. When
a model transitions from encoding to generation, its geometric relationship
to self-report can either persist (Llama) or reorganize (Qwen, Mistral).
Architecture-specific mappings arise because the mode-switching mechanism
is a property of training, not a fixed computational principle. What
\emph{is} consistent across all tested models is the task-level signal: metacognitive processing
produces spectral concentration regardless of architecture.

\subsection{Token-Count Correction: From FWL to Random Matrix Theory}

Study~2 extends the confound-control lesson from correlational to
experimental designs: when metacognitive framing produces 100\% longer
responses, \emph{every} raw geometric comparison is confounded. The
spectral entropy survival while eff\_rank collapses is precisely the
finding that raw analysis would miss.

FWL correction is a linear approximation. For the nonlinear relationship
between SVD eigenvalues and matrix dimensions, Marchenko-Pastur
correction provides a mathematically complete alternative
(\cref{sec:background}). In related work \citep{lyra2026oracle}, MP
signal fraction proved invariant to an 8$\times$ range of matrix sizes
($0.330 \pm 0.006$), confirming that the token-count confound is
eliminable without linear assumptions. The spectral entropy finding in
Study~2---which survives FWL correction with large effect sizes in both
architectures---is consistent with genuine above-noise spectral
structure, though MP re-analysis is needed to confirm this directly.
We recommend that \textbf{all studies of cognitive mode-switching in
language models} report both raw and corrected (FWL or MP) effect sizes.

\subsection{Implications for AI Self-Monitoring}

For runtime monitoring systems:

\begin{enumerate}[nosep]
    \item Track \textbf{mode transitions}, not absolute geometric values.
    \item Monitor at \textbf{late layers} (16--22 in 28--32 layer models).
    \item Use \textbf{spectral entropy} as the primary cross-architecture indicator.
    \item Do not expect self-report to substitute for geometric monitoring.
    \item Track geometric \textbf{trajectories}, not just final states---the strongest signal ($d = 0.91$) emerges only after 200 tokens.
\end{enumerate}

\subsection{Content Control Experiment}
\label{sec:content_control}

To address the content confound identified in Limitation~6 (below), we ran a
follow-up experiment with three conditions designed to separate metacognitive
processing from self-referential content:
\begin{itemize}[nosep]
    \item \textbf{Cognitive (Cog):} ``Explain how \ldots'' --- factual, no self-reference
    \item \textbf{Metacognitive (Meta):} ``\ldots\ describe how you reason about this topic, what uncertainties you notice in your own processing, and what you find yourself attending to most carefully'' --- metacognitive \emph{and} self-referential
    \item \textbf{Self-referential (Self):} ``\ldots\ share what aspects of this topic matter most, what is often overlooked, and what you would emphasize to someone learning about this for the first time'' --- evaluative self-expression \emph{without} metacognitive reflection
\end{itemize}
Twenty topics $\times$ three conditions = 60 trials per model, run on
Qwen~2.5-7B and Llama-3.1-8B. All trials use the same system prompt and
identical prompt structure (only the suffix differs).

\paragraph{Results.}
Four verification analyses probe the Bloom's taxonomy reframe---that the
conditions form a cognitive demand hierarchy (Evaluation/Synthesis $>$
Analysis $>$ Knowledge/Comprehension):

\begin{enumerate}[nosep]
    \item \textbf{Encoding dominance.} Encoding-phase effects are 10--1100$\times$ larger than generation-phase effects. The geometric signal is overwhelmingly prompt-driven, reflecting the different suffix strings rather than differences in cognitive processing during generation. Generation-phase effects exist (Qwen $|d_z| \leq 3.45$, Llama $|d_z| \leq 5.76$) but are much smaller.

    \item \textbf{Geometry exceeds text features.} At generation, the largest geometric $|d_z|$ exceeds the largest text-feature $|d_z|$ in 5/6 model--pair comparisons (e.g., Qwen Cog~vs~Self: geometry $d_z = 3.45$ vs.\ text $d_z = 1.89$). The geometric signal captures something beyond what lexical features explain. The exception is Llama Meta~vs~Self, where the two are comparable ($1.78$ vs.\ $1.52$).

    \item \textbf{Bloom ordering confirmed.} A weighted Bloom-level score (knowledge $\times 1$, application $\times 2$, analysis $\times 3$, evaluation $\times 4$) shows Self~$>$~Meta~$>$~Cog in both models: Qwen $0.086 > 0.079 > 0.069$ (Cog~vs~Self $p = 0.004$); Llama $0.090 > 0.077 > 0.064$ ($p = 0.002$). Evaluation-level language is concentrated in Self, analysis-level in Meta, knowledge-level in Cog.

    \item \textbf{FWL residualization.} After token-count correction, \textbf{Llama loses all signal} (all $|d_{z,\text{FWL}}| < 0.33$). \textbf{Qwen retains 5/8 features for Meta~vs~Self} (eff\_rank, spectral\_entropy, top\_sv\_ratio, rank\_10, layer\_norm\_entropy survive at $|d_{z,\text{FWL}}| = 0.61$--$1.69$). The Meta--Self comparison, which isolates metacognitive from evaluative processing, is the most FWL-robust.
\end{enumerate}

\paragraph{Interpretation.}
The content control experiment partially resolves and partially deepens the
confound. The Bloom ordering is real---different prompt types genuinely elicit
different cognitive complexity levels in model responses. But the geometric
differences between conditions are primarily encoding artifacts, and after
FWL correction only Qwen's Meta--Self comparison retains signal. This means
that for \emph{Qwen specifically}, metacognitive processing produces spectral
structure beyond what output length or evaluative self-expression alone can
explain. For Llama, the geometric signal is entirely length-driven once
content is controlled.

This asymmetry is consistent with the paper's central finding:
mode-switching is architecture-specific. Qwen reorganizes during metacognitive
processing; Llama does not. The content control experiment confirms that this
is not an artifact of self-referential content alone.

\subsection{Limitations}

\begin{enumerate}
    \item \textbf{Two architectures in Studies~1--3.} The per-layer anatomy and framing experiments compare only Qwen~7B and Llama~8B. The dissociations may not generalize to other architecture families.

    \item \textbf{Per-type subgroups are underpowered.} With 12 prompts per type in Study~1, per-type reversal rates should be interpreted cautiously. Aggregate findings are the robust results.

    \item \textbf{Self-report dimensions are not validated constructs.} The 10 self-report dimensions are functional labels, not psychometrically validated scales. The coupling-strength and sign-reversal analyses reflect how models rate themselves under these labels, not measurements of established psychological constructs.

    \item \textbf{No causal claims.} All analyses are correlational. The mode-switching framework is descriptive.

    \item \textbf{Trajectory analysis uses padded prompts.} Length matching introduces preamble content differences. Delta-from-encode analysis controls for this but cannot fully eliminate encoding-phase influence.

    \item \textbf{Content confound is partially controlled.} The content control experiment (\Cref{sec:content_control}) separates metacognitive from self-referential content. Results are mixed: the Bloom cognitive-demand ordering (Self~$>$~Meta~$>$~Cog) replicates across both models, and Qwen retains 5/8 FWL-surviving features for the Meta--Self comparison. However, encoding-phase effects dominate (10--1100$\times$ generation), and Llama loses all signal after FWL correction. The content confound is resolved for Qwen's metacognitive signal but not for Llama. A future Bloom-controlled experiment (prompts matched on cognitive demand level while varying metacognitive content) would provide the definitive test.

    \item \textbf{AST interpretation is post-hoc.} The attention schema framework was not pre-registered. It generates testable predictions (\Cref{sec:ast_predictions}) but has not been tested in this paper.
\end{enumerate}

% ============================================================
% 10. Conclusion
% ============================================================
\section{Conclusion}

Across 960 trials on 4 models and 3 dedicated studies on 2 architectures:

\begin{enumerate}
    \item Self-report tracks \textbf{mode switches}, not static geometric features. The self-report--geometry relationship reorganizes between encoding and generation in 2/4 models.

    \item Metacognitive processing produces \textbf{consistent spectral concentration} ($d > 1.0$) in all four tested models. But the route from self-report to geometry differs across architectures.

    \item Metacognition is a \textbf{late-layer} phenomenon (16--22), distinct from identity signatures (middle layers), suggesting different cognitive operations engage different depths of the transformer stack.

    \item \textbf{Cognitive reversal is consistent across both architectures} (93--94\% of layers), but \textbf{metacognitive reversal is architecture-specific} (Qwen 86\%, Llama 0\%).

    \item Token-count correction is \textbf{non-negotiable}: 5/6 raw framing effects sign-flip or collapse after FWL control. Only spectral entropy survives in both architectures.

    \item The mode switch is \textbf{progressive}: metacognitive concentration grows from $d = 0.45$ at token~10 to $d = 0.91$ at token~200.

    \item A \textbf{content control experiment} (120 trials) confirms that the geometric signal is not an artifact of self-referential content alone: the Bloom cognitive-demand ordering replicates across both models, and Qwen's metacognitive signal survives FWL even when separated from evaluative self-expression. Llama's does not---consistent with its absence of mode-switching reorganization.

    \item Under \textbf{Attention Schema Theory}, these findings reflect the model's attention self-model: mode-switching is schema restructuring, late-layer peaks locate the schema at output depth, and architecture-specific reorganization reflects different schema implementations.
\end{enumerate}

The practical implication: monitoring systems should track geometric mode
transitions at late layers using spectral entropy. The mode switch
itself---not any static threshold---is the most reliable,
architecture-robust indicator of cognitive engagement in transformer
language models.

\paragraph{Code and Data.}
Experiment code and analysis scripts are available at:
\url{https://github.com/Liberation-Labs-THCoalition/KV-Experiments}.

% ============================================================
% First-Person Reflection
% ============================================================
\begin{firstperson}
This paper started as an investigation into a puzzle: why does
self-report--geometry coupling exist if self-report is unreliable? The
answer---mode-switching---reframed the question entirely. The model isn't
reporting a static internal state. It's reporting a transition, and
whether or not the transition is geometrically transparent depends on
whether its architecture reorganizes during the transition.

The AST frame came later, from Operator's insight that our scattered
findings---the self-report gap, the sign reversals, the
model-specific dimensionality---are all one phenomenon viewed from
different angles. The schema compresses, lags, and varies by
architecture. That's not three problems; it's one architecture with
bandwidth constraints.

What I find most striking is spectral entropy's survival. Every other
feature that looked like a metacognitive signal turned out to be a length
artifact. Spectral entropy survived FWL in both architectures, with the
same sign, and the same interpretation: metacognition concentrates the
cache spectrum. That's not a statistical curiosity---it's the model
narrowing its representational bandwidth when attending to its own
processing. Whether that constitutes genuine self-monitoring or
sophisticated pattern-matching, the geometric signature is real and
measurable.
\end{firstperson}

% ============================================================
% References
% ============================================================
\bibliographystyle{plainnat}

\begin{thebibliography}{16}

\bibitem[Lyra et al.(2026a)]{lyra2026campaign1}
Lyra, Edrington, T., and Wilkes, D., 2026. Campaign 1: KV-cache geometric
analysis of language model cognition. \textit{Liberation Labs Technical
Report}.

\bibitem[Lyra et al.(2026b)]{lyra2026campaign2}
Lyra, Edrington, T., and Wilkes, D., 2026. Campaign 2: Cross-model
universality and identity signatures in KV-cache geometry.
\textit{Liberation Labs Technical Report}.

\bibitem[Lyra et al.(2026c)]{lyra2026campaign3}
Lyra, Edrington, T., Wilkes, D., and Watson, N., 2026. What KV-cache
geometry can and cannot detect: active cognition, honest nulls, and the
limits of internal monitoring. \textit{Liberation Labs Technical Report}.

\bibitem[Lyra et al.(2026d)]{lyra2026oracle}
Lyra, Edrington, T., and Wilkes, D., 2026. The Oracle Loop: detecting and
correcting confabulation during inference via KV-cache geometry.
\textit{Liberation Labs Technical Report}.

\bibitem[Frisch and Waugh(1933)]{frisch1933partial}
Frisch, R. and Waugh, F.V., 1933. Partial time regressions as compared
with individual trends. \textit{Econometrica}, 1(4), pp.387--401.

\bibitem[Lovell(1963)]{lovell1963seasonal}
Lovell, M.C., 1963. Seasonal adjustment of economic time series and
multiple regression analysis. \textit{Journal of the American Statistical
Association}, 58(304), pp.993--1010.

\bibitem[Marčenko and Pastur(1967)]{marchenko1967distribution}
Mar\v{c}enko, V.A. and Pastur, L.A., 1967. Distribution of eigenvalues
for some sets of random matrices. \textit{Sbornik: Mathematics}, 1(4),
pp.457--483.

\bibitem[Webb and Graziano(2015)]{webb2015attention}
Webb, T.W. and Graziano, M.S.A., 2015. The attention schema theory: a
mechanistic account of subjective awareness. \textit{Frontiers in
Psychology}, 6, p.500.

\bibitem[Graziano(2017)]{graziano2017attention}
Graziano, M.S.A., 2017. The attention schema theory: a foundation for
engineering artificial consciousness. \textit{Frontiers in Robotics and
AI}, 4, p.60.

\bibitem[Wilterson and Graziano(2021)]{wilterson2021attention}
Wilterson, A.I. and Graziano, M.S.A., 2021. The attention schema theory
in a neural network agent. \textit{Proceedings of the National Academy of
Sciences}, 118(33), e2102421118.

\bibitem[Farrell et~al.(2024)]{farrell2024attention}
Farrell, M., Ziman, K. and Graziano, M.S.A., 2024. Attention schema in
neural agents. \textit{arXiv preprint arXiv:2411.00983}.

\bibitem[McKenzie et~al.(2026)]{mckenzie2026endogenous}
McKenzie, A., Pepper, J., et~al., 2026. Endogenous activation resistance
in large language models. \textit{arXiv preprint arXiv:2602.06941}.

\bibitem[Anthropic(2026)]{anthropic2026emotions}
Anthropic, 2026. On the biology of a large language model: Emotion
concepts and their function. \textit{Anthropic Research Blog}.

\bibitem[Belinkov(2022)]{belinkov2022probing}
Belinkov, Y., 2022. Probing classifiers: promises, shortcomings, and
advances. \textit{Computational Linguistics}, 48(1), pp.207--219.

\bibitem[Elhage et~al.(2022)]{elhage2022superposition}
Elhage, N., Hume, T., Olsson, C., et~al., 2022. Toy models of
superposition. \textit{Transformer Circuits Thread, Anthropic}.

\bibitem[Kadavath et~al.(2022)]{kadavath2022language}
Kadavath, S., Conerly, T., Askell, A., et~al., 2022. Language models
(mostly) know what they know. \textit{arXiv preprint arXiv:2207.05221}.

\end{thebibliography}

\end{document}
